\chapter{Definizione del servizio}
\label{chap:servizio}

\section{Dominio applicativo}
\label{sec:dominio}

Il servizio si colloca nel dominio \emph{Finance \& Enterprise Document Compliance},
con particolare riferimento al sotto-settore della \emph{procure-to-pay automation}
per imprese di medie dimensioni. Il sistema si propone come piattaforma SaaS multi-tenant
per la ricezione, conservazione, classificazione e ricerca di documenti amministrativi,
con un livello di \emph{intelligent triage} integrato.

I documenti gestiti appartengono tipicamente alle seguenti categorie:
\begin{itemize}[leftmargin=*]
    \item \textbf{Documenti fiscali}: fatture attive e passive, note di credito, estratti conto.
    \item \textbf{Documenti contrattuali}: contratti con fornitori, NDA, ordini di acquisto.
    \item \textbf{Comunicazioni di compliance}: circolari normative, relazioni di audit,
          certificazioni.
    \item \textbf{Documentazione operativa}: verbali di riunione, report periodici,
          documentazione tecnica interna.
\end{itemize}

\section{Stakeholder}
\label{sec:stakeholder}

\begin{description}[leftmargin=*]
    \item[Utente finale (knowledge worker)] addetto amministrativo, contabile o
          responsabile di funzione che riceve via e-mail fatture, contratti e
          comunicazioni di compliance e necessita di una classificazione e
          archiviazione tempestiva. Interagisce con il sistema tramite interfaccia web
          per caricare documenti, consultare l'archivio, condividere risorse e revisionare
          le proposte dell'agente.

    \item[Domain administrator] referente IT del dominio aziendale, responsabile
          della configurazione delle utenze, della governance dei branding e
          della consultazione dei \emph{metrics} aggregati a livello di tenant.
          Il primo utente a registrarsi con un determinato dominio email assume
          automaticamente questo ruolo (\emph{first-user-wins}).

    \item[Auditor / Compliance officer] consulta il \emph{trail} delle
          operazioni dell'agente e degli utenti per verifiche periodiche di correttezza
          della classificazione e tracciabilit\`a delle azioni.

    \item[Operatore di piattaforma] presidia la salute del sistema attraverso
          gli \emph{health endpoint} e i log strutturati prodotti dai worker.
          \TODO{Con l'agente di monitoring Kubernetes, questa figura interagir\`a
          anche con le proposte di remediation del secondo agente}
\end{description}

\section{Workload e principali eventi}
\label{sec:workload}

Il sistema gestisce le seguenti classi di eventi, con i relativi profili di carico attesi:

\begin{itemize}[leftmargin=*]
    \item \textbf{Sincronizzazione IMAP}: il \emph{poller} interroga periodicamente
          gli account configurati. Per utenti attivi (\texttt{last\_active\_at} recente)
          la frequenza \`e adattiva; per utenti inattivi la cadenza scende a 24~ore per
          contenere i costi di rete e di IMAP API usage.

    \item \textbf{Job di classificazione}: ogni e-mail con allegati genera un record
          \texttt{email\_jobs} che fluisce attraverso la macchina a stati
          \texttt{pending} $\rightarrow$ \texttt{processing} $\rightarrow$
          \texttt{done}/\texttt{failed}.

    \item \textbf{Caricamento manuale diretto}: l'utente seleziona un file dal browser.
          Il backend genera un presigned URL verso SeaweedFS; il browser carica direttamente
          il file sull'object storage; infine il backend viene notificato per finalizzare il
          record nel database. Questo pattern two-phase elimina il transito del binario
          attraverso il backend.

    \item \textbf{Caricamento con classificazione AI}: l'utente usa l'endpoint
          \texttt{/documents/upload-ai}. Il backend restituisce immediatamente un
          \texttt{202 Accepted} e crea un record \texttt{agent\_files} in stato
          \texttt{pending} che sar\`a processato asincronamente dall'agent worker.

    \item \textbf{Operazioni utente su metadati}: creazione e spostamento di cartelle,
          download di documenti, gestione dei permessi di condivisione, revisione delle
          \emph{proposals} dell'agente nell'\emph{inbox view}.

    \item \textbf{Ricerca}: l'utente invia query testuali che vengono eseguite su
          OpenSearch con combinazione di \emph{full-text} e ricerca vettoriale k-NN.

    \item \textbf{Operazioni amministrative}: gestione utenze del dominio,
          consultazione metrics e audit log.
\end{itemize}

La distribuzione attesa del carico \`e fortemente \emph{bursty} (picchi a fine mese
e a fine trimestre fiscale) e prevalentemente asincrona, condizione favorevole
all'adozione di un'architettura a code persistenti basata su polling del database.

\section{Scenari di guasto principali}
\label{sec:scenari_guasto}

Sono stati individuati i seguenti scenari critici, analizzati in dettaglio nel
\Cref{chap:failure}:

\begin{enumerate}[leftmargin=*]
    \item \emph{Indisponibilit\`a del database principale}: blocco delle
          operazioni di scrittura sui metadati; il health endpoint espone lo stato
          degradato con HTTP 503.
    \item \emph{Indisponibilit\`a dell'object storage}: impossibilit\`a di
          archiviare nuovi allegati o servire download; i metadati rimangono
          accessibili.
    \item \emph{Errore o congestione del provider LLM}: latenza elevata,
          \emph{rate limit}, errori transienti che impediscono la classificazione.
          Il sistema degrada a modalit\`a ``manual filing''.
    \item \emph{Allegato malformato}: PDF corrotti, archivi non standard, documenti
          scansionati senza testo estraibile. L'agente abbassa la confidenza e
          invia il documento in inbox.
    \item \emph{Hallucination dell'agente}: classificazione semanticamente plausibile
          ma errata. Il meccanismo di validazione UUID e il confidence threshold
          costituiscono la prima e seconda linea di difesa.
    \item \emph{Fault del worker}: crash durante l'elaborazione di un job.
          Il job resta in stato \texttt{processing} fino a intervento manuale.
    \item \emph{Server IMAP irraggiungibile o credenziali revocate}: il poller
          registra l'errore, aggiorna \texttt{last\_synced\_at} e prosegue con
          gli altri account.
\end{enumerate}

\section{Perch\'e l'affidabilit\`a \`e critica}
\label{sec:affidabilita_critica}

Nel dominio considerato un singolo errore di archiviazione pu\`o tradursi in:
\emph{(i)} ritardo di pagamento di una fattura con conseguente penale o sospensione
di fornitura; \emph{(ii)} smarrimento di documenti contrattuali con perdita di valore
probatorio; \emph{(iii)} violazioni della normativa di conservazione (e.g.\ regolamenti
fiscali nazionali e GDPR per i dati personali contenuti nei documenti).

La piattaforma agisce dunque come \emph{system of record} e deve garantire
\textbf{durabilit\`a}, \textbf{tracciabilit\`a} e \textbf{disponibilit\`a} elevate,
anche in presenza di degrado del livello agentico.
